\chapter{State Eilenberg Theory for Residual Mealy Recognition}
\label{ch:state-eilenberg-residual-mealy-recognition}

\section{Overview}
\label{sec:state-eilenberg-overview}

The preceding chapters constructed, for internal categories
\[
\mathbb A
\qquad\text{and}\qquad
\mathbb B
\]
in a Grothendieck topos \(\mathcal E\), the final behaviour system
\[
\mathsf{Beh}_{\mathbb A,\mathbb B}
\]
for Mealy morphisms
\[
\mathbb A\rightsquigarrow\mathbb B.
\]
They also associated to each behaviour specification
\[
S\hookrightarrow \mathsf{Beh}_{\mathbb A,\mathbb B}
\]
its canonical residual closure
\[
\mathsf D(S)\hookrightarrow \mathsf{Beh}_{\mathbb A,\mathbb B}.
\]
This residual closure is the canonical separated state recognizer generated
by the specified behaviours. It is separated because it lives as a subobject
of the final behaviour system, and it is residual because it is closed under
all derivatives induced by input arrows of \(\mathbb A\).

The notation \(\mathsf D(S)\) denotes the same residual closure denoted
\(D(S)\) in the previous chapter. In this chapter we use the typography
\(\mathsf D\) to emphasize its role as a closure operator on behaviour
specifications.

This chapter develops the Eilenberg theory at the state-recognizer level.
There are two layers.

The first layer consists of behaviour specifications
\[
S\hookrightarrow \mathsf{Beh}_{\mathbb A,\mathbb B}.
\]
These are the observable behaviour objects whose closure properties one
wishes to study.

The second layer consists of closed residual state recognizers
\[
D\hookrightarrow \mathsf{Beh}_{\mathbb A,\mathbb B}.
\]
A specification \(S\) determines such a recognizer by residual closure:
\[
\mathsf D(S)\hookrightarrow \mathsf{Beh}_{\mathbb A,\mathbb B}.
\]
Its associated residual syntactic state category is
\[
\operatorname{SynState}_{\mathbb A,\mathbb B}(S)
:=
\int_{\mathbb A}\mathsf D(S).
\]
This is a state/process category. Its objects are residual states, and its
arrows remember both a residual state and an input arrow of \(\mathbb A\)
which acts on that state.

Thus the state-level Eilenberg picture is
\[
\text{behaviour specifications}
\quad\longleftrightarrow\quad
\text{closed residual state recognizers}.
\]
The correspondence is mediated by the residual closure operator
\[
S\longmapsto \mathsf D(S).
\]

The main theorem of this chapter states that, relative to a chosen class of
state sizes and a chosen doctrine of behaviour operations, stable classes of
closed residual state recognizers correspond exactly to recognizable
behaviour varieties. More precisely, if
\[
\mathcal K_Z\subseteq \mathcal E/Z,
\qquad
Z=A_0\times B_0,
\]
is an appropriate state-size class, and if \(\mathcal O\) is a doctrine of
behaviour operations, then there is an order isomorphism
\[
\left\{
\begin{array}{c}
\text{stable classes of closed}\\
\mathcal K\text{-state recognizers}
\end{array}
\right\}
\cong
\left\{
\begin{array}{c}
\mathcal K\text{-recognizable}\\
\mathcal O\text{-Mealy behaviour varieties}
\end{array}
\right\}.
\]
The assignment from recognizers to behaviours is
\[
\mathcal P
\longmapsto
\mathcal V_{\mathcal P}
:=
\left\{
S\hookrightarrow\mathsf{Beh}_{\mathbb A,\mathbb B}
\ \middle|\
\mathsf D(S)\in\mathcal P
\right\},
\]
and the assignment from behaviours to recognizers is
\[
\mathcal V
\longmapsto
\mathcal P_{\mathcal V}
:=
\left\{
D\in\mathsf{SRec}^{\mathcal K}_{\mathbb A,\mathbb B}
\ \middle|\
D\in\mathcal V
\right\}.
\]

The quotient of input arrows by their induced transition-output operation is
not studied in this chapter. That quotient-level construction belongs to the
Reiterman theory of the next chapter. The present chapter stops at the
state-recognizer level:
\[
S
\quad\rightsquigarrow\quad
\mathsf D(S)
\quad\rightsquigarrow\quad
\int_{\mathbb A}\mathsf D(S).
\]
The next chapter will pass from a closed residual state recognizer \(D\) to
the effective transition-output quotient determined by the input-labelled
action category
\[
\int_{\mathbb A}D\longrightarrow \mathbb A^{\mathrm{op}}.
\]
In the one-object case this specializes to the quotient of the input monoid
by the syntactic transition-output congruence.

\section{Ambient setting and residual notation}
\label{sec:state-eilenberg-ambient-setting}

Throughout this chapter, assume that
\[
\mathcal E
\]
is a Grothendieck topos. Thus \(\mathcal E\) is locally cartesian closed,
regular, Barr-exact, and locally presentable. Every slice
\[
\mathcal E/Z
\]
is again regular and Barr-exact.

We work in a fixed universe. All indexed products, intersections, and
families appearing below are taken over small indexing sets in this universe.
A collection means a small collection unless explicitly stated otherwise.

Let
\[
\mathbb A
=
(A_0,A_1,s_A,t_A,e_A,m_A)
\qquad\text{and}\qquad
\mathbb B
=
(B_0,B_1,s_B,t_B,e_B,m_B)
\]
be internal categories in \(\mathcal E\). We use the composition convention
from the preceding chapters: if
\[
a:u\to v,
\qquad
b:v\to w,
\]
then
\[
b\circ a:u\to w
\]
means first \(a\), then \(b\). Equivalently,
\[
m_A(a,b)=b\circ a.
\]

Write
\[
Z:=A_0\times B_0.
\]
Let
\[
\mathsf{Beh}
:=
\mathsf{Beh}_{\mathbb A,\mathbb B}
\]
denote the final behaviour system for Mealy morphisms
\[
\mathbb A\rightsquigarrow\mathbb B.
\]
Its carrier is an object of
\[
\mathcal E/Z.
\]
We write its projections as
\[
p_{\mathsf{Beh}}:\mathsf{Beh}\to A_0,
\qquad
q_{\mathsf{Beh}}:\mathsf{Beh}\to B_0.
\]

The transition and output maps of \(\mathsf{Beh}\) are
\[
\delta_{\mathsf{Beh}}:
A_1\times_{t_A,A_0,p_{\mathsf{Beh}}}\mathsf{Beh}
\to
\mathsf{Beh}
\]
and
\[
\omega_{\mathsf{Beh}}:
A_1\times_{t_A,A_0,p_{\mathsf{Beh}}}\mathsf{Beh}
\to
B_1.
\]
For an admissible input arrow
\[
a:u\to v
\]
and a behaviour \(\beta\) lying over \(v\), write
\[
\partial_a\beta
:=
\delta_{\mathsf{Beh}}(a,\beta).
\]
The derivative \(\partial_a\beta\) lies over \(u\).

The derivative laws are
\[
\partial_{e_A(v)}\beta=\beta
\]
and, for composable arrows
\[
a:u\to v,
\qquad
b:v\to w,
\]
and \(\beta\) lying over \(w\),
\[
\partial_{b\circ a}\beta
=
\partial_a(\partial_b\beta).
\]

Let
\[
S\hookrightarrow\mathsf{Beh}
\]
be a behaviour specification in
\[
\mathcal E/Z.
\]
Write
\[
p_S:S\to A_0
\]
for the composite with \(p_{\mathsf{Beh}}\). Define
\[
A_1\ltimes S
:=
A_1\times_{t_A,A_0,p_S}S.
\]
A generalized element of \(A_1\ltimes S\) is a pair
\[
(a,\beta),
\]
where
\[
a:u\to v
\]
is an arrow of \(\mathbb A\), and \(\beta\in S\) lies over \(v\).

The derivative evaluation map is
\[
d_S:A_1\ltimes S\to\mathsf{Beh},
\qquad
d_S(a,\beta)=\partial_a\beta.
\]
The object \(A_1\ltimes S\) is regarded as an object of
\[
\mathcal E/Z
\]
through the composite
\[
A_1\ltimes S
\xrightarrow{d_S}
\mathsf{Beh}
\to Z.
\]
With this slice structure, \(d_S\) is a morphism in \(\mathcal E/Z\).

\begin{definition}[Residual closure]
\label{def:state-eilenberg-residual-closure}
The \textbf{residual closure} of a behaviour specification
\[
S\hookrightarrow\mathsf{Beh}
\]
is
\[
\mathsf D(S)
:=
\operatorname{im}
\left(
A_1\ltimes S
\xrightarrow{d_S}
\mathsf{Beh}
\right)
\]
in \(\mathcal E/Z\). Thus there is a regular epi--mono factorization
\[
A_1\ltimes S
\twoheadrightarrow
\mathsf D(S)
\hookrightarrow
\mathsf{Beh}.
\]
\end{definition}

The assignment
\[
S\longmapsto\mathsf D(S)
\]
is a closure operator on
\[
\operatorname{Sub}_Z(\mathsf{Beh}).
\]
Thus, for subobjects \(S,T\hookrightarrow\mathsf{Beh}\), one has:
\[
S\leq \mathsf D(S),
\]
\[
S\leq T
\quad\Longrightarrow\quad
\mathsf D(S)\leq \mathsf D(T),
\]
and
\[
\mathsf D(\mathsf D(S))=\mathsf D(S).
\]
The closed subobjects are precisely the sub-Mealy systems of
\[
\mathsf{Beh}.
\]

Equivalently, the inclusion of closed behaviour systems into
\[
\operatorname{Sub}_Z(\mathsf{Beh})
\]
has left adjoint \(\mathsf D\). Thus, for every closed behaviour system
\[
D\hookrightarrow\mathsf{Beh},
\]
one has
\[
\mathsf D(S)\leq D
\quad\Longleftrightarrow\quad
S\leq D.
\]

\begin{definition}[Closed behaviour system]
\label{def:state-eilenberg-closed-behaviour-system}
A subobject
\[
D\hookrightarrow \mathsf{Beh}
\]
in \(\mathcal E/Z\) is called a \textbf{closed behaviour system} if
\[
\mathsf D(D)=D.
\]
Equivalently, \(D\) is a sub-Mealy system of the final behaviour system.
\end{definition}

Thus a closed behaviour system is a behaviourally separated residual state
recognizer. It has no distinct states with the same final behaviour, because
its carrier is already a subobject of the final behaviour object.

\section{Recognition by Mealy systems}
\label{sec:state-eilenberg-recognition-by-mealy-systems}

Let
\[
X:\mathbb A\rightsquigarrow\mathbb B
\]
be a Mealy system. Its behaviour map is the unique Mealy transformation
\[
\operatorname{beh}_X:X\to\mathsf{Beh}.
\]

\begin{definition}[Realized behaviour object]
\label{def:state-eilenberg-realized-behaviour-object}
The \textbf{realized behaviour object} of \(X\) is the image
\[
\operatorname{Real}(X)
:=
\operatorname{im}
\left(
X\xrightarrow{\operatorname{beh}_X}\mathsf{Beh}
\right)
\]
in \(\mathcal E/Z\). Thus there is a regular epi--mono factorization
\[
X\twoheadrightarrow
\operatorname{Real}(X)
\hookrightarrow
\mathsf{Beh}.
\]
\end{definition}

The object \(\operatorname{Real}(X)\) is the separated behavioural image of
the system \(X\). Passing from \(X\) to \(\operatorname{Real}(X)\) identifies
states of \(X\) with the same final behaviour.

\begin{definition}[Recognition by a Mealy system]
\label{def:state-eilenberg-recognition}
A Mealy system
\[
X:\mathbb A\rightsquigarrow\mathbb B
\]
\textbf{recognizes} a behaviour specification
\[
S\hookrightarrow\mathsf{Beh}
\]
if
\[
S\leq\operatorname{Real}(X)
\]
as subobjects of \(\mathsf{Beh}\) in \(\mathcal E/Z\).
\end{definition}

Thus recognition is inclusion of the specified behaviours into the separated
image of the recognizing system.

\begin{proposition}[Realized behaviours are closed]
\label{prop:state-eilenberg-realized-behaviours-closed}
For every Mealy system
\[
X:\mathbb A\rightsquigarrow\mathbb B,
\]
the subobject
\[
\operatorname{Real}(X)\hookrightarrow\mathsf{Beh}
\]
is a closed behaviour system. Equivalently,
\[
\mathsf D(\operatorname{Real}(X))
=
\operatorname{Real}(X).
\]
\end{proposition}

\begin{proof}
Factor the behaviour map as
\[
X\twoheadrightarrow I\hookrightarrow\mathsf{Beh},
\]
where
\[
I=\operatorname{Real}(X).
\]
Since regular epimorphisms are stable under pullback, the induced map
\[
A_1\times_{t_A,A_0,p_X}X
\twoheadrightarrow
A_1\times_{t_A,A_0,p_I}I
\]
is a regular epimorphism.

The derivative map of \(\mathsf{Beh}\), after pulling back along this regular
epimorphism, agrees with the composite
\[
A_1\times_{t_A,A_0,p_X}X
\xrightarrow{\delta_X}
X
\twoheadrightarrow
I.
\]
This is precisely the transition-preservation equation for the Mealy
transformation
\[
\operatorname{beh}_X:X\to\mathsf{Beh}.
\]
Therefore, after pullback along a regular epimorphism, the derivative of a
state of \(I\) lands in \(I\).

Pullback along a regular epimorphism reflects containment of subobjects in a
regular category. Hence the derivative map of \(\mathsf{Beh}\) restricts to
\[
A_1\times_{t_A,A_0,p_I}I\to I.
\]
The output map restricts from the output map of \(\mathsf{Beh}\), and the
Mealy equations hold because they hold in \(\mathsf{Beh}\). Thus \(I\) is a
sub-Mealy system of \(\mathsf{Beh}\). Equivalently,
\[
\mathsf D(I)=I.
\]
\end{proof}

\begin{corollary}
\label{cor:state-eilenberg-recognition-implies-closure-contained}
If a Mealy system
\[
X:\mathbb A\rightsquigarrow\mathbb B
\]
recognizes
\[
S\hookrightarrow\mathsf{Beh},
\]
then
\[
\mathsf D(S)\leq \operatorname{Real}(X).
\]
\end{corollary}

\begin{proof}
Since \(X\) recognizes \(S\), one has
\[
S\leq\operatorname{Real}(X).
\]
By Proposition~\ref{prop:state-eilenberg-realized-behaviours-closed},
\[
\operatorname{Real}(X)
\]
is closed. Since
\[
\mathsf D(S)
\]
is the smallest closed behaviour system containing \(S\), it follows that
\[
\mathsf D(S)\leq\operatorname{Real}(X).
\]
\end{proof}

\section{State recognizers and recognizable specifications}
\label{sec:state-eilenberg-state-recognizers}

The preceding section shows that every Mealy system recognizes specifications
through a closed subobject of the final behaviour system. We now abstract this
closed subobject as the state-recognizer level of the theory.

\begin{definition}[State-size class]
\label{def:state-eilenberg-state-size-class}
A \textbf{state-size class} over
\[
Z=A_0\times B_0
\]
is a full subcollection
\[
\mathcal K_Z\subseteq\mathcal E/Z
\]
closed under isomorphism. Additional closure properties are specified when
needed.
\end{definition}

The class \(\mathcal K_Z\) records the allowed size or structure of state
objects. Examples include finite objects in a Set-based situation, compact
objects, presentable objects, or other internally specified classes in a
topos.

The basic state-recognizability criterion below only requires closure of
\[
\mathcal K_Z
\]
under subobjects, because recognizers are already taken to be closed
subobjects of the final behaviour system. If one also wants to pass from
arbitrary \(\mathcal K\)-state Mealy systems to their separated behavioural
images, then \(\mathcal K_Z\) should additionally be closed under regular
images.

\begin{definition}[\(\mathcal K\)-state recognizer]
\label{def:state-eilenberg-K-state-recognizer}
A \textbf{\(\mathcal K\)-state recognizer} for
\[
(\mathbb A,\mathbb B)
\]
is a closed behaviour subobject
\[
D\hookrightarrow\mathsf{Beh}_{\mathbb A,\mathbb B}
\]
whose carrier lies in
\[
\mathcal K_Z.
\]
The poset of such recognizers is denoted
\[
\mathsf{SRec}^{\mathcal K}_{\mathbb A,\mathbb B}.
\]
\end{definition}

Thus a \(\mathcal K\)-state recognizer is already behaviourally separated:
it is not an arbitrary Mealy system, but a sub-Mealy system of the final
behaviour system.

\begin{definition}[Recognition by a state recognizer]
\label{def:state-eilenberg-recognition-by-state-recognizer}
A \(\mathcal K\)-state recognizer
\[
D\hookrightarrow\mathsf{Beh}
\]
\textbf{recognizes} a behaviour specification
\[
S\hookrightarrow\mathsf{Beh}
\]
if
\[
S\leq D.
\]
\end{definition}

\begin{proposition}[State recognizability criterion]
\label{prop:state-eilenberg-state-recognizability-criterion}
Assume that
\[
\mathcal K_Z
\]
is closed under subobjects. A behaviour specification
\[
S\hookrightarrow\mathsf{Beh}
\]
is recognized by a \(\mathcal K\)-state recognizer if and only if
\[
\mathsf D(S)\in\mathcal K_Z.
\]
\end{proposition}

\begin{proof}
If
\[
\mathsf D(S)\in\mathcal K_Z,
\]
then \(\mathsf D(S)\) is itself a closed behaviour subobject containing
\(S\), hence a \(\mathcal K\)-state recognizer recognizing \(S\).

Conversely, suppose that \(S\) is recognized by a \(\mathcal K\)-state
recognizer \(D\). Then
\[
S\leq D.
\]
Since \(D\) is closed, minimality of residual closure gives
\[
\mathsf D(S)\leq D.
\]
Because \(\mathcal K_Z\) is closed under subobjects and \(D\in\mathcal K_Z\),
it follows that
\[
\mathsf D(S)\in\mathcal K_Z.
\]
\end{proof}

The criterion says that recognizability is completely controlled by the
residual closure of the specification. In particular, any non-separated
recognizer can be replaced, for recognition purposes, by its separated
behavioural image. If the allowed state-size class is also closed under
regular images, then this replacement preserves membership in the allowed
state-size class.

\section{State divisibility and residual syntactic state categories}
\label{sec:state-eilenberg-state-divisibility}

The residual closure
\[
\mathsf D(S)
\]
is not merely a recognizer of \(S\). It is the canonical minimal separated
residual recognizer of \(S\). This minimality can be expressed as a
divisibility statement.

\begin{definition}[State divisibility]
\label{def:state-eilenberg-state-divisibility}
Let
\[
Y
\qquad\text{and}\qquad
X
\]
be Mealy systems from
\[
\mathbb A
\]
to
\[
\mathbb B.
\]
We say that
\[
Y
\]
\textbf{state-divides}
\[
X
\]
if there exists a sub-Mealy system
\[
X'\hookrightarrow X
\]
and a regular epimorphism of Mealy systems
\[
X'\twoheadrightarrow Y.
\]
\end{definition}

\begin{theorem}[State syntactic divisibility]
\label{thm:state-eilenberg-state-syntactic-divisibility}
Let
\[
S\hookrightarrow\mathsf{Beh}
\]
be a behaviour specification. If a Mealy system
\[
X:\mathbb A\rightsquigarrow\mathbb B
\]
recognizes \(S\), then the canonical residual state recognizer
\[
\mathsf D(S)
\]
state-divides \(X\).
\end{theorem}

\begin{proof}
By Corollary~\ref{cor:state-eilenberg-recognition-implies-closure-contained},
\[
\mathsf D(S)\leq\operatorname{Real}(X).
\]
Write the image factorization of the behaviour map as
\[
X\twoheadrightarrow I\hookrightarrow \mathsf{Beh},
\qquad
I=\operatorname{Real}(X).
\]
Form the pullback in \(\mathcal E/Z\):
\[
X_S
:=
X\times_{\mathsf{Beh}}\mathsf D(S).
\]
Since
\[
\mathsf D(S)\leq I\leq\mathsf{Beh},
\]
this pullback identifies with
\[
X_S\cong X\times_I \mathsf D(S).
\]
The map
\[
X\twoheadrightarrow I
\]
is a regular epimorphism, and regular epimorphisms are stable under pullback
in the slice \(\mathcal E/Z\). Hence the projection
\[
X_S\to \mathsf D(S)
\]
is a regular epimorphism.

Since
\[
\mathsf D(S)\hookrightarrow\mathsf{Beh}
\]
is a sub-Mealy system and
\[
\operatorname{beh}_X:X\to\mathsf{Beh}
\]
is a Mealy transformation, the pullback \(X_S\) inherits a sub-Mealy-system
structure from \(X\). Thus
\[
X_S\hookrightarrow X
\]
is a sub-Mealy system, and
\[
X_S\twoheadrightarrow \mathsf D(S)
\]
is a regular epimorphism of Mealy systems. Therefore \(\mathsf D(S)\)
state-divides \(X\).
\end{proof}

\begin{definition}[Residual syntactic state category]
\label{def:state-eilenberg-synstate}
The \textbf{residual syntactic state category} of a specification
\[
S\hookrightarrow\mathsf{Beh}
\]
is
\[
\operatorname{SynState}_{\mathbb A,\mathbb B}(S)
:=
\int_{\mathbb A}\mathsf D(S).
\]
It is the action category of the residual state system
\[
\mathsf D(S).
\]
\end{definition}

Explicitly, the objects of
\[
\operatorname{SynState}_{\mathbb A,\mathbb B}(S)
\]
are residual states
\[
\beta\in \mathsf D(S).
\]
An arrow is an admissible pair
\[
(a,\beta),
\]
where
\[
a:u\to v
\]
is an arrow of \(\mathbb A\), and \(\beta\) lies over \(v\). Its source is
\[
\beta,
\]
and its target is
\[
\partial_a\beta.
\]
The output labelling is inherited from the output map of the final behaviour
system. The canonical input projection is a functor
\[
\int_{\mathbb A}\mathsf D(S)\longrightarrow \mathbb A^{\mathrm{op}},
\]
since the arrow \(a:u\to v\) of \(\mathbb A\) is read as an arrow
\(v\to u\) in \(\mathbb A^{\mathrm{op}}\).

\begin{corollary}[Labelled state divisibility]
\label{cor:state-eilenberg-labelled-state-divisibility}
If \(X\) recognizes \(S\), then
\[
\operatorname{SynState}_{\mathbb A,\mathbb B}(S)
\]
is obtained from
\[
\int_{\mathbb A}X
\]
by passing to a labelled internal subcategory and then to a labelled regular
quotient.
\end{corollary}

\begin{proof}
Apply the action-category construction
\[
\int_{\mathbb A}(-)
\]
to the sub-Mealy system
\[
X_S\hookrightarrow X
\]
and the regular epimorphism
\[
X_S\twoheadrightarrow\mathsf D(S)
\]
from Theorem~\ref{thm:state-eilenberg-state-syntactic-divisibility}. This
gives
\[
\int_{\mathbb A}X_S
\hookrightarrow
\int_{\mathbb A}X
\]
and
\[
\int_{\mathbb A}X_S
\twoheadrightarrow
\int_{\mathbb A}\mathsf D(S).
\]
The first map is a labelled internal subcategory inclusion.

For the second map, the object map is the regular epimorphism
\[
X_S\twoheadrightarrow\mathsf D(S).
\]
The arrow map is obtained by pulling this regular epimorphism back along
\[
t_A:A_1\to A_0,
\]
hence is again a regular epimorphism. Therefore the induced labelled functor
is a levelwise regular quotient. It preserves the input and output labellings
because the quotient map is a Mealy transformation. Since
\[
\int_{\mathbb A}\mathsf D(S)
=
\operatorname{SynState}_{\mathbb A,\mathbb B}(S),
\]
the claim follows.
\end{proof}

\section{Behaviour-operation doctrines}
\label{sec:state-eilenberg-behaviour-operation-doctrines}

The Eilenberg correspondence is formulated relative to a specified family of
behaviour operations. This section packages those operations and their
recognizer-level counterparts as a single formal structure.

\begin{definition}[Behaviour-operation doctrine]
\label{def:state-eilenberg-behaviour-operation-doctrine}
Fix
\[
(\mathbb A,\mathbb B)
\]
and write
\[
\mathsf{Beh}:=\mathsf{Beh}_{\mathbb A,\mathbb B}.
\]
A \textbf{behaviour-operation doctrine}
\[
\mathcal O
\]
consists of:
\begin{enumerate}
\item a small set of operation symbols \(\alpha\);
\item for each operation symbol \(\alpha\), a small arity set \(I_\alpha\);
\item for each \(\alpha\), a monotone operation on behaviour subobjects
\[
\Omega_\alpha:
\prod_{i\in I_\alpha}
\operatorname{Sub}_Z(\mathsf{Beh})
\to
\operatorname{Sub}_Z(\mathsf{Beh}).
\]
\end{enumerate}
Given closed behaviour systems
\[
(D_i\hookrightarrow\mathsf{Beh})_{i\in I_\alpha},
\]
the associated state-recognizer constructor is defined by
\[
\Omega_\alpha^{\mathrm{st}}
\bigl((D_i)_{i\in I_\alpha}\bigr)
:=
\mathsf D
\left(
\Omega_\alpha
\bigl((D_i)_{i\in I_\alpha}\bigr)
\right).
\]
\end{definition}

Thus the state-level operation is not an additional primitive structure. It
is the closed residual realization of the behaviour-level operation applied
to closed recognizers.

\begin{remark}
\label{rem:state-eilenberg-operations-need-not-preserve-closed}
The operation
\[
\Omega_\alpha
\]
need not preserve closed behaviour systems. The state operation
\[
\Omega_\alpha^{\mathrm{st}}
\]
is therefore defined by applying the behaviour operation first and then
reflecting back into closed behaviour systems by residual closure.
\end{remark}

\begin{definition}[\(\mathcal O\)-closed state-size class]
\label{def:state-eilenberg-O-closed-state-size-class}
Let
\[
\mathcal K_Z\subseteq\mathcal E/Z
\]
be a state-size class closed under subobjects. We say that
\[
\mathcal K_Z
\]
is \textbf{\(\mathcal O\)-closed} if, for every operation symbol \(\alpha\)
and every family of closed \(\mathcal K\)-state recognizers
\[
(D_i)_{i\in I_\alpha},
\]
the carrier of
\[
\Omega_\alpha^{\mathrm{st}}
\bigl((D_i)_{i\in I_\alpha}\bigr)
\]
again lies in
\[
\mathcal K_Z.
\]
\end{definition}

The closure requirement says that the chosen notion of recognizable state
object is stable under the residual closure of the specified behaviour
operations. This is the state-recognizer analogue of closure of a language
variety under Boolean operations, derivatives, inverse images, or other
chosen behaviour-forming operations in classical settings.

\section{Local behaviour varieties}
\label{sec:state-eilenberg-local-behaviour-varieties}

Fix a behaviour-operation doctrine
\[
\mathcal O
\]
and an \(\mathcal O\)-closed state-size class
\[
\mathcal K_Z\subseteq \mathcal E/Z.
\]

\begin{definition}[Local behaviour variety]
\label{def:state-eilenberg-local-behaviour-variety}
A \textbf{local \(\mathcal K\)-recognizable \(\mathcal O\)-Mealy behaviour
variety} is a class
\[
\mathcal V_{\mathbb A,\mathbb B}
\subseteq
\operatorname{Sub}_Z(\mathsf{Beh}_{\mathbb A,\mathbb B})
\]
satisfying the following conditions.

\begin{enumerate}
\item \textbf{Downward closure.} If
\[
S\in\mathcal V_{\mathbb A,\mathbb B}
\]
and
\[
T\leq S,
\]
then
\[
T\in\mathcal V_{\mathbb A,\mathbb B}.
\]

\item \textbf{Residual closure.} If
\[
S\in\mathcal V_{\mathbb A,\mathbb B},
\]
then
\[
\mathsf D(S)\in\mathcal V_{\mathbb A,\mathbb B}.
\]

\item \textbf{\(\mathcal O\)-operation closure.} For every operation symbol
\(\alpha\), if
\[
(S_i)_{i\in I_\alpha}
\]
is a family of elements of \(\mathcal V_{\mathbb A,\mathbb B}\), then
\[
\Omega_\alpha\bigl((S_i)_{i\in I_\alpha}\bigr)
\in
\mathcal V_{\mathbb A,\mathbb B}.
\]

\item \textbf{\(\mathcal K\)-recognizability.} For every
\[
S\in\mathcal V_{\mathbb A,\mathbb B},
\]
one has
\[
\mathsf D(S)\in\mathcal K_Z.
\]
\end{enumerate}
\end{definition}

The elements of a behaviour variety are specifications, not recognizers.
Recognizability is imposed by requiring the residual closure
\[
\mathsf D(S)
\]
to have carrier in the allowed state-size class
\[
\mathcal K_Z.
\]

\section{Stable state-recognizer classes}
\label{sec:state-eilenberg-stable-state-recognizer-classes}

The recognizer-side object corresponding to a behaviour variety is a stable
class of closed state recognizers.

\begin{definition}[Stable state-recognizer class]
\label{def:state-eilenberg-stable-state-recognizer-class}
A \textbf{local stable \(\mathcal K\)-state-recognizer class} is a class
\[
\mathcal P_{\mathbb A,\mathbb B}
\subseteq
\mathsf{SRec}^{\mathcal K}_{\mathbb A,\mathbb B}
\]
closed under:
\begin{enumerate}
\item isomorphism, if representatives of subobjects are being used;
\item closed subrecognizers;
\item the state operations
\[
\Omega_\alpha^{\mathrm{st}}
\]
for all operation symbols \(\alpha\in\mathcal O\).
\end{enumerate}
\end{definition}

Here a \textbf{closed subrecognizer} of
\[
D\hookrightarrow\mathsf{Beh}
\]
is a closed behaviour system
\[
E\hookrightarrow\mathsf{Beh}
\]
with
\[
E\leq D.
\]

Given a local stable state-recognizer class
\[
\mathcal P_{\mathbb A,\mathbb B},
\]
define
\[
\mathcal V_{\mathcal P,\mathbb A,\mathbb B}
:=
\left\{
S\leq\mathsf{Beh}_{\mathbb A,\mathbb B}
\ \middle|\
\mathsf D(S)\in\mathcal P_{\mathbb A,\mathbb B}
\right\}.
\]
Conversely, given a local behaviour variety
\[
\mathcal V_{\mathbb A,\mathbb B},
\]
define
\[
\mathcal P_{\mathcal V,\mathbb A,\mathbb B}
:=
\left\{
D\in\mathsf{SRec}^{\mathcal K}_{\mathbb A,\mathbb B}
\ \middle|\
D\in\mathcal V_{\mathbb A,\mathbb B}
\right\}.
\]

\section{The local state Eilenberg theorem}
\label{sec:state-eilenberg-local-state-correspondence}

\begin{theorem}[Local state Eilenberg theorem]
\label{thm:state-eilenberg-state-correspondence}
Assume that the state-size class
\[
\mathcal K_Z
\]
is closed under subobjects and is \(\mathcal O\)-closed. Then the assignments
\[
\mathcal P\longmapsto\mathcal V_{\mathcal P}
\]
and
\[
\mathcal V\longmapsto\mathcal P_{\mathcal V}
\]
define mutually inverse order isomorphisms, ordered by inclusion, between
local stable \(\mathcal K\)-state-recognizer classes and local
\(\mathcal K\)-recognizable \(\mathcal O\)-Mealy behaviour varieties.
\end{theorem}

\begin{proof}
Let
\[
\mathcal P_{\mathbb A,\mathbb B}
\subseteq
\mathsf{SRec}^{\mathcal K}_{\mathbb A,\mathbb B}
\]
be a local stable state-recognizer class. We first show that
\[
\mathcal V_{\mathcal P,\mathbb A,\mathbb B}
\]
is a local behaviour variety.

If
\[
S\in\mathcal V_{\mathcal P}
\]
and
\[
T\leq S,
\]
then by monotonicity of residual closure,
\[
\mathsf D(T)\leq\mathsf D(S).
\]
Since
\[
\mathsf D(S)\in\mathcal P,
\]
and since \(\mathsf D(T)\) is a closed subrecognizer of \(\mathsf D(S)\),
closure of \(\mathcal P\) under closed subrecognizers gives
\[
\mathsf D(T)\in\mathcal P.
\]
Here \(\mathsf D(T)\) lies in \(\mathcal K_Z\) because \(\mathsf D(T)\) is a
subobject of \(\mathsf D(S)\) and \(\mathcal K_Z\) is closed under
subobjects. Hence
\[
T\in\mathcal V_{\mathcal P}.
\]

Residual closure follows from idempotence:
if
\[
S\in\mathcal V_{\mathcal P},
\]
then
\[
\mathsf D(S)\in\mathcal P.
\]
But
\[
\mathsf D(\mathsf D(S))=\mathsf D(S),
\]
so
\[
\mathsf D(S)\in\mathcal V_{\mathcal P}.
\]

For \(\mathcal O\)-operation closure, let
\[
(S_i)_{i\in I_\alpha}
\]
be a family in
\[
\mathcal V_{\mathcal P}.
\]
Then
\[
\mathsf D(S_i)\in\mathcal P
\]
for every \(i\). Since
\[
S_i\leq\mathsf D(S_i),
\]
monotonicity of \(\Omega_\alpha\) gives
\[
\Omega_\alpha((S_i)_{i\in I_\alpha})
\leq
\Omega_\alpha((\mathsf D(S_i))_{i\in I_\alpha}).
\]
Applying monotonicity of \(\mathsf D\), we get
\[
\mathsf D\bigl(\Omega_\alpha((S_i)_{i\in I_\alpha})\bigr)
\leq
\mathsf D\bigl(\Omega_\alpha((\mathsf D(S_i))_{i\in I_\alpha})\bigr).
\]
By definition of the state operation, the right-hand side is
\[
\Omega_\alpha^{\mathrm{st}}
\bigl((\mathsf D(S_i))_{i\in I_\alpha}\bigr).
\]
This belongs to \(\mathcal P\) because \(\mathcal P\) is closed under the
state operation \(\Omega_\alpha^{\mathrm{st}}\). Since the left-hand side is
a closed subrecognizer of the right-hand side, and since \(\mathcal K_Z\) is
closed under subobjects, closure of \(\mathcal P\) under closed
subrecognizers gives
\[
\mathsf D\bigl(\Omega_\alpha((S_i)_{i\in I_\alpha})\bigr)
\in\mathcal P.
\]
Hence
\[
\Omega_\alpha((S_i)_{i\in I_\alpha})\in\mathcal V_{\mathcal P}.
\]

Finally,
\[
S\in\mathcal V_{\mathcal P}
\]
implies
\[
\mathsf D(S)\in\mathcal P\subseteq
\mathsf{SRec}^{\mathcal K},
\]
so \(S\) is \(\mathcal K\)-recognizable.

Conversely, let
\[
\mathcal V_{\mathbb A,\mathbb B}
\]
be a local behaviour variety. We show that
\[
\mathcal P_{\mathcal V,\mathbb A,\mathbb B}
\]
is a stable state-recognizer class.

Closure under isomorphism is immediate if representatives of subobjects are
used. If
\[
E\leq D,
\]
where \(D\in\mathcal P_{\mathcal V}\) and \(E\) is a closed
\(\mathcal K\)-state recognizer, then
\[
D\in\mathcal V
\]
and downward closure gives
\[
E\in\mathcal V.
\]
Hence
\[
E\in\mathcal P_{\mathcal V}.
\]

For operation closure, let
\[
(D_i)_{i\in I_\alpha}
\]
be a family in \(\mathcal P_{\mathcal V}\). Then each \(D_i\) belongs to
\[
\mathcal V.
\]
By \(\mathcal O\)-operation closure,
\[
\Omega_\alpha((D_i)_{i\in I_\alpha})\in\mathcal V.
\]
By residual closure of \(\mathcal V\),
\[
\mathsf D(\Omega_\alpha((D_i)_{i\in I_\alpha}))
=
\Omega_\alpha^{\mathrm{st}}
\bigl((D_i)_{i\in I_\alpha}\bigr)
\]
also belongs to \(\mathcal V\). Since \(\mathcal K_Z\) is
\(\mathcal O\)-closed, its carrier lies in \(\mathcal K_Z\). Thus
\[
\Omega_\alpha^{\mathrm{st}}
\bigl((D_i)_{i\in I_\alpha}\bigr)
\in\mathcal P_{\mathcal V}.
\]

It remains to show that the assignments are inverse. Let \(\mathcal V\) be a
behaviour variety. Then
\[
S\in\mathcal V_{\mathcal P_{\mathcal V}}
\]
if and only if
\[
\mathsf D(S)\in\mathcal P_{\mathcal V},
\]
if and only if
\[
\mathsf D(S)\in\mathcal V.
\]
Since \(\mathcal V\) is residual closed and downward closed, and since
\[
S\leq\mathsf D(S),
\]
we have
\[
\mathsf D(S)\in\mathcal V
\quad\Longleftrightarrow\quad
S\in\mathcal V.
\]
Therefore
\[
\mathcal V_{\mathcal P_{\mathcal V}}=\mathcal V.
\]

Now let \(\mathcal P\) be a stable state-recognizer class. For a closed
\(\mathcal K\)-state recognizer \(D\),
\[
D\in\mathcal P_{\mathcal V_{\mathcal P}}
\]
if and only if
\[
D\in\mathcal V_{\mathcal P},
\]
if and only if
\[
\mathsf D(D)\in\mathcal P.
\]
Since \(D\) is closed,
\[
\mathsf D(D)=D.
\]
Hence
\[
D\in\mathcal P_{\mathcal V_{\mathcal P}}
\quad\Longleftrightarrow\quad
D\in\mathcal P.
\]
Thus
\[
\mathcal P_{\mathcal V_{\mathcal P}}=\mathcal P.
\]

Both assignments are monotone with respect to inclusion of classes, and the
displayed inverse equalities show that they define an order isomorphism.
\end{proof}

\section{Global interfaces and state-admissible translations}
\label{sec:state-eilenberg-global-interfaces-state-translations}

We now pass from the local theory at a fixed pair
\[
(\mathbb A,\mathbb B)
\]
to a global theory over varying Mealy interfaces.

\begin{definition}[Mealy interface]
\label{def:state-eilenberg-mealy-interface}
A \textbf{Mealy interface} is a pair
\[
\Xi=(\mathbb A_\Xi,\mathbb B_\Xi)
\]
of internal categories in the ambient topos. We write
\[
\mathsf{Beh}_\Xi
:=
\mathsf{Beh}_{\mathbb A_\Xi,\mathbb B_\Xi},
\]
\[
Z_\Xi:=
(A_\Xi)_0\times(B_\Xi)_0,
\]
and
\[
\mathsf D_\Xi
\]
for the residual closure operator on
\[
\operatorname{Sub}_{Z_\Xi}(\mathsf{Beh}_\Xi).
\]
\end{definition}

\begin{definition}[Specification poset of an interface]
\label{def:state-eilenberg-specification-poset-interface}
For an interface \(\Xi\), define
\[
\mathsf{Spec}_\Xi
:=
\operatorname{Sub}_{Z_\Xi}(\mathsf{Beh}_\Xi).
\]
This is the poset of behaviour specifications at \(\Xi\).
\end{definition}

\begin{definition}[State-admissible interface translation]
\label{def:state-eilenberg-state-admissible-interface-translation}
Let
\[
\Xi'=(\mathbb A',\mathbb B')
\qquad\text{and}\qquad
\Xi=(\mathbb A,\mathbb B)
\]
be interfaces. A \textbf{state-admissible interface translation}
\[
\Gamma:\Xi'\to\Xi
\]
consists of a monotone reindexing operation on behaviour specifications
\[
\Gamma_{\mathrm{beh}}^\ast:
\mathsf{Spec}_{\Xi}
\to
\mathsf{Spec}_{\Xi'}
\]
satisfying the following axioms.

\begin{enumerate}
\item \textbf{Compatibility with residual closure.} For every
\[
S\in\mathsf{Spec}_{\Xi},
\]
one has
\[
\Gamma_{\mathrm{beh}}^\ast(\mathsf D_\Xi(S))
=
\mathsf D_{\Xi'}(\Gamma_{\mathrm{beh}}^\ast S).
\]

\item \textbf{Compatibility with behaviour operations.} For every operation
symbol \(\alpha\),
\[
\Gamma_{\mathrm{beh}}^\ast
\left(
\Omega_\alpha^\Xi((S_i)_{i\in I_\alpha})
\right)
=
\Omega_\alpha^{\Xi'}
\left(
(\Gamma_{\mathrm{beh}}^\ast S_i)_{i\in I_\alpha}
\right).
\]
\end{enumerate}
\end{definition}

\begin{remark}[Variance]
\label{rem:state-eilenberg-admissible-translation-variance}
The operation
\[
\Gamma_{\mathrm{beh}}^\ast
\]
is specified directly on behaviour subobjects. It need not be literal
pullback along a displayed map of final behaviour objects. In concrete
situations it may arise from pullback, image, restriction, precomposition, or
another behaviour-reindexing construction. The axioms above record precisely
the variance and compatibility required for the global state Eilenberg
theory.
\end{remark}

If
\[
D\hookrightarrow\mathsf{Beh}_\Xi
\]
is a closed behaviour system, define its reindexed closed recognizer by
\[
\Gamma_{\mathrm{st}}^\ast D
:=
\mathsf D_{\Xi'}(\Gamma_{\mathrm{beh}}^\ast D).
\]
By residual compatibility, if \(D\) is closed then
\[
\Gamma_{\mathrm{st}}^\ast D
=
\Gamma_{\mathrm{beh}}^\ast D.
\]
Indeed,
\[
\Gamma_{\mathrm{beh}}^\ast D
=
\Gamma_{\mathrm{beh}}^\ast(\mathsf D_\Xi D)
=
\mathsf D_{\Xi'}(\Gamma_{\mathrm{beh}}^\ast D).
\]

\begin{definition}[Functorial system of state-admissible translations]
\label{def:state-eilenberg-functorial-system-translations}
A collection of state-admissible translations is called \textbf{functorial}
if it is closed under identities and composition and satisfies
\[
(\mathrm{id}_\Xi)_{\mathrm{beh}}^\ast
=
\mathrm{id}_{\mathsf{Spec}_\Xi}
\]
and, for composable translations
\[
\Xi''\xrightarrow{\Gamma'}\Xi'\xrightarrow{\Gamma}\Xi,
\]
one has
\[
(\Gamma\circ\Gamma')_{\mathrm{beh}}^\ast
=
(\Gamma')_{\mathrm{beh}}^\ast\circ \Gamma_{\mathrm{beh}}^\ast.
\]
\end{definition}

\section{Global state Eilenberg correspondence}
\label{sec:state-eilenberg-global-state-correspondence}

Let
\[
\mathfrak T
\]
be a category of Mealy interfaces and state-admissible interface
translations. Thus the objects of \(\mathfrak T\) are interfaces \(\Xi\), and
the morphisms are state-admissible translations
\[
\Gamma:\Xi'\to\Xi
\]
forming a functorial system in the sense of
Definition~\ref{def:state-eilenberg-functorial-system-translations}.

For each interface \(\Xi\), suppose we are given a state-size class
\[
\mathcal K_{Z_\Xi}\subseteq\mathcal E/Z_\Xi
\]
closed under subobjects, and a behaviour-operation doctrine
\[
\mathcal O_\Xi.
\]
Assume that the doctrines are compatible with state-admissible translations
as in
Definition~\ref{def:state-eilenberg-state-admissible-interface-translation}.

We also assume that translations preserve the chosen state-size classes on
closed recognizers: if
\[
D\in\mathsf{SRec}^{\mathcal K}_\Xi,
\]
then
\[
\Gamma_{\mathrm{st}}^\ast D
\in
\mathsf{SRec}^{\mathcal K}_{\Xi'}.
\]
Equivalently, the carrier of \(\Gamma_{\mathrm{st}}^\ast D\) lies in
\[
\mathcal K_{Z_{\Xi'}}.
\]

\begin{definition}[Global behaviour variety]
\label{def:state-eilenberg-global-behaviour-variety}
A \textbf{global \(\mathcal K\)-recognizable \(\mathcal O\)-Mealy behaviour
variety} on \(\mathfrak T\) assigns to every interface \(\Xi\) a local
\(\mathcal K\)-recognizable \(\mathcal O_\Xi\)-Mealy behaviour variety
\[
\mathcal V_\Xi
\subseteq
\mathsf{Spec}_\Xi
\]
such that for every state-admissible translation
\[
\Gamma:\Xi'\to\Xi
\]
and every
\[
S\in\mathcal V_\Xi,
\]
one has
\[
\Gamma_{\mathrm{beh}}^\ast S\in\mathcal V_{\Xi'}.
\]
\end{definition}

\begin{definition}[Global stable state-recognizer class]
\label{def:state-eilenberg-global-stable-state-recognizer-class}
A \textbf{global stable \(\mathcal K\)-state-recognizer class} on
\(\mathfrak T\) assigns to every interface \(\Xi\) a local stable
\(\mathcal K\)-state-recognizer class
\[
\mathcal P_\Xi
\subseteq
\mathsf{SRec}^{\mathcal K}_\Xi
\]
such that for every state-admissible translation
\[
\Gamma:\Xi'\to\Xi
\]
and every
\[
D\in\mathcal P_\Xi,
\]
one has
\[
\Gamma_{\mathrm{st}}^\ast D\in\mathcal P_{\Xi'}.
\]
\end{definition}

Given a global stable state-recognizer class
\[
\mathcal P=(\mathcal P_\Xi)_\Xi,
\]
define
\[
\mathcal V_{\mathcal P,\Xi}
:=
\left\{
S\in\mathsf{Spec}_\Xi
\ \middle|\
\mathsf D_\Xi(S)\in\mathcal P_\Xi
\right\}.
\]
Conversely, given a global behaviour variety
\[
\mathcal V=(\mathcal V_\Xi)_\Xi,
\]
define
\[
\mathcal P_{\mathcal V,\Xi}
:=
\left\{
D\in\mathsf{SRec}^{\mathcal K}_\Xi
\ \middle|\
D\in\mathcal V_\Xi
\right\}.
\]

\begin{theorem}[Global state Eilenberg theorem]
\label{thm:state-eilenberg-global-state-correspondence}
Assume that, for every interface \(\Xi\), the relevant state-size class
\[
\mathcal K_{Z_\Xi}
\]
is closed under subobjects and is \(\mathcal O_\Xi\)-closed. Assume also
that state-admissible translations preserve the chosen state-size classes on
closed recognizers. Then the assignments
\[
\mathcal P\longmapsto\mathcal V_{\mathcal P}
\]
and
\[
\mathcal V\longmapsto\mathcal P_{\mathcal V}
\]
define mutually inverse order isomorphisms, ordered objectwise by inclusion,
between global stable \(\mathcal K\)-state-recognizer classes and global
\(\mathcal K\)-recognizable \(\mathcal O\)-Mealy behaviour varieties.
\end{theorem}

\begin{proof}
Objectwise, this is the local correspondence of
Theorem~\ref{thm:state-eilenberg-state-correspondence}. It remains only to
verify the inverse-image conditions.

Let
\[
\mathcal P
\]
be a global stable state-recognizer class and suppose
\[
S\in\mathcal V_{\mathcal P,\Xi}.
\]
Then
\[
\mathsf D_\Xi(S)\in\mathcal P_\Xi.
\]
For a state-admissible translation
\[
\Gamma:\Xi'\to\Xi,
\]
we compute
\[
\mathsf D_{\Xi'}(\Gamma_{\mathrm{beh}}^\ast S)
=
\Gamma_{\mathrm{beh}}^\ast(\mathsf D_\Xi(S))
\]
by residual compatibility. Since \(\mathsf D_\Xi(S)\) is closed, this is the
same closed recognizer as
\[
\Gamma_{\mathrm{st}}^\ast(\mathsf D_\Xi(S)).
\]
By global stability of \(\mathcal P\), it belongs to \(\mathcal P_{\Xi'}\).
Hence
\[
\Gamma_{\mathrm{beh}}^\ast S
\in
\mathcal V_{\mathcal P,\Xi'}.
\]
Thus \(\mathcal V_{\mathcal P}\) is global.

Conversely, let
\[
\mathcal V
\]
be a global behaviour variety and suppose
\[
D\in\mathcal P_{\mathcal V,\Xi}.
\]
Then
\[
D\in\mathcal V_\Xi.
\]
For a state-admissible translation
\[
\Gamma:\Xi'\to\Xi,
\]
global inverse-image stability gives
\[
\Gamma_{\mathrm{beh}}^\ast D\in\mathcal V_{\Xi'}.
\]
By residual closure of \(\mathcal V_{\Xi'}\),
\[
\Gamma_{\mathrm{st}}^\ast D
=
\mathsf D_{\Xi'}(\Gamma_{\mathrm{beh}}^\ast D)
\in
\mathcal V_{\Xi'}.
\]
By the assumed preservation of the chosen state-size classes on closed
recognizers,
\[
\Gamma_{\mathrm{st}}^\ast D
\in
\mathsf{SRec}^{\mathcal K}_{\Xi'}.
\]
Therefore
\[
\Gamma_{\mathrm{st}}^\ast D\in\mathcal P_{\mathcal V,\Xi'}.
\]
Thus \(\mathcal P_{\mathcal V}\) is global.

The fact that the assignments are inverse follows objectwise from
Theorem~\ref{thm:state-eilenberg-state-correspondence}. The order
isomorphism is also objectwise.
\end{proof}

\section{The one-object Set specialization}
\label{sec:state-eilenberg-one-object-set}

Let
\[
\mathcal E=\mathbf{Set},
\qquad
\mathbb A=M,
\qquad
\mathbb B=N,
\]
where \(M\) and \(N\) are monoids, regarded as one-object internal
categories. We retain the convention
\[
mn=m\circ n.
\]
Thus \(mn\) means first \(n\), then \(m\) as categorical composition, while
the right-action execution by \(mn\) is computed as first \(m\), then \(n\).

In this one-object case, a Mealy system has an ordinary state set \(X\), a
transition function
\[
\delta:M\times X\to X,
\]
and an output function
\[
\omega:M\times X\to N,
\]
satisfying the Mealy unit and multiplication laws
\[
\delta(1,x)=x,
\qquad
\omega(1,x)=1,
\]
and
\[
\delta(mn,x)=\delta(n,\delta(m,x)),
\]
\[
\omega(mn,x)=\omega(m,x)\omega(n,\delta(m,x)).
\]
Equivalently, writing
\[
x\cdot m:=\delta(m,x),
\]
the transition is a right action in the operational order:
\[
x\cdot(mn)=(x\cdot m)\cdot n.
\]

A behaviour in the final behaviour system can be described as an
\(N\)-valued residual cocycle
\[
\lambda:M\times M\to N
\]
satisfying
\[
\lambda(r,1)=1,
\]
and
\[
\lambda(r,mn)=\lambda(r,m)\lambda(rm,n).
\]
Here \(r\) records the residual position and \(m\) records the next input to
be processed.

Derivatives are given by
\[
(\partial_m\lambda)(r,n)=\lambda(mr,n).
\]
More generally,
\[
(\partial_r\lambda)(s,m)=\lambda(rs,m).
\]
Consequently, for
\[
S\subseteq\mathsf{Beh}_{M,N},
\]
the residual closure is
\[
\mathsf D(S)
=
\{\partial_r\lambda\mid r\in M,\ \lambda\in S\}.
\]

At the residual state
\[
\partial_r\lambda,
\]
processing an input
\[
m\in M
\]
gives
\[
\partial_r\lambda\cdot m
=
\partial_{rm}\lambda,
\]
and the emitted output is
\[
\omega(m,\partial_r\lambda)
=
\lambda(r,m).
\]

In this specialization, the state Eilenberg theorem says that recognizable
classes of \(N\)-valued Mealy behaviours are controlled by their residual
closures inside the final behaviour set. A specification
\[
S\subseteq\mathsf{Beh}_{M,N}
\]
is recognized by a state recognizer of the allowed size class precisely when
the residual state set
\[
\mathsf D(S)
\]
belongs to that size class.

The residual syntactic state category
\[
\operatorname{SynState}_{M,N}(S)
=
\int_M \mathsf D(S)
\]
is the action category of the residual right action of \(M\) on
\[
\mathsf D(S).
\]
It records states
\[
\lambda\in\mathsf D(S)
\]
and transitions
\[
\lambda
\longrightarrow
\lambda\cdot m
\]
labelled by inputs \(m\in M\), with output determined by the corresponding
residual cocycle.

\section{The free-input specialization}
\label{sec:state-eilenberg-free-input-specialization}

Suppose now that
\[
M=I^\ast
\]
is the free monoid on a set \(I\). Then the final behaviour object has the
generator-level normal form
\[
\mathsf{Beh}_{I^\ast,N}
\cong
N^{I^+}.
\]
Here \(I^+\) denotes the set of nonempty words.

A behaviour is a function
\[
\beta:I^+\to N.
\]
Word concatenation is written in operational order: \(wu\) means first process
\(w\), then process \(u\). The derivative by a word
\[
w\in I^\ast
\]
is the residual shift
\[
(\partial_w\beta)(u)=\beta(wu),
\qquad
u\in I^+.
\]

For
\[
S\subseteq N^{I^+},
\]
the residual closure is
\[
\mathsf D(S)
=
\{\partial_w\beta\mid w\in I^\ast,\ \beta\in S\}.
\]

Thus the residual syntactic state category is
\[
\operatorname{SynState}_{I^\ast,N}(S)
=
\int_{I^\ast}\mathsf D(S).
\]
Its objects are residual behaviours
\[
\partial_w\beta,
\]
and its arrows are generated by shifts along input words. In particular, the
generator-level transitions are
\[
\partial_w\beta
\longrightarrow
\partial_{wi}\beta
\]
for \(i\in I\).

For a generator
\[
i\in I,
\]
the output at state \(\partial_w\beta\) is
\[
\omega(i,\partial_w\beta)=\beta(wi).
\]
For a composite word
\[
u=i_1\cdots i_n,
\]
the corresponding word-output is the coherent cocycle value
\[
\lambda_\beta(w,u)
=
\beta(wi_1)\beta(wi_1i_2)\cdots
\beta(wi_1\cdots i_n).
\]
Thus the value \(\beta(wu)\) records the generator-level output at the end of
the residual prefix \(wu\), while the folded output along the whole word
\(u\) is the product of the successive generator-level outputs.

In the classical finite-state situation, if \(I\), \(N\), and
\[
\mathsf D(S)
\]
are finite, then \(\mathsf D(S)\) is the usual finite residual state set
generated by the specified behaviours. The Eilenberg theorem of this chapter
abstracts precisely this residual-state construction to internal categories
and arbitrary state-size doctrines.

\section{The state Eilenberg triangle}
\label{sec:state-eilenberg-triangle}

The state-level part of the Eilenberg--Reiterman picture has two layers:
\[
\text{behaviour specifications}
\quad\longleftrightarrow\quad
\text{residual state recognizers}.
\]
The first layer consists of subobjects
\[
S\hookrightarrow\mathsf{Beh}_{\mathbb A,\mathbb B}.
\]
The second layer consists of closed behaviour systems
\[
D\hookrightarrow\mathsf{Beh}_{\mathbb A,\mathbb B}.
\]

The bridge between them is the residual closure operator
\[
S\longmapsto \mathsf D(S).
\]
A specification is recognized by an allowed state recognizer exactly when its
residual closure has allowed state size.

The syntactic state recognizer of \(S\) is
\[
\mathsf D(S),
\]
and the corresponding residual syntactic state category is
\[
\operatorname{SynState}_{\mathbb A,\mathbb B}(S)
=
\int_{\mathbb A}\mathsf D(S).
\]
This category remembers residual states and their input-driven transitions.

The Eilenberg correspondence classifies stable classes of such closed state
recognizers. It does not yet identify input arrows which induce the same
operation on all residual states. That quotient-level identification is the
subject of the next chapter.

\section{Main consolidation}
\label{sec:state-eilenberg-main-consolidation}

\begin{theorem}[State Eilenberg theory for residual Mealy recognition]
\label{thm:state-eilenberg-main-consolidation}
Let
\[
\mathcal E
\]
be a Grothendieck topos, and let
\[
\mathbb A,\mathbb B
\]
be internal categories in \(\mathcal E\). Let
\[
\mathsf{Beh}_{\mathbb A,\mathbb B}
\]
be the final behaviour system for Mealy morphisms
\[
\mathbb A\rightsquigarrow\mathbb B.
\]
Then:

\begin{enumerate}
\item Every Mealy system
\[
X:\mathbb A\rightsquigarrow\mathbb B
\]
has a realized behaviour subobject
\[
\operatorname{Real}(X)
\hookrightarrow
\mathsf{Beh}_{\mathbb A,\mathbb B}.
\]

\item The realized behaviour object is a closed sub-Mealy system of the final
behaviour system:
\[
\mathsf D(\operatorname{Real}(X))
=
\operatorname{Real}(X).
\]

\item A Mealy system \(X\) recognizes a specification
\[
S\hookrightarrow\mathsf{Beh}_{\mathbb A,\mathbb B}
\]
if and only if
\[
S\leq\operatorname{Real}(X).
\]

\item If \(X\) recognizes \(S\), then
\[
\mathsf D(S)\leq\operatorname{Real}(X).
\]

\item If \(X\) recognizes \(S\), then
\[
\mathsf D(S)
\]
is a regular quotient of a sub-Mealy system of \(X\). Equivalently,
\[
\mathsf D(S)
\]
state-divides every Mealy system recognizing \(S\).

\item The residual syntactic state category of \(S\) is
\[
\operatorname{SynState}_{\mathbb A,\mathbb B}(S)
=
\int_{\mathbb A}\mathsf D(S).
\]

\item If \(X\) recognizes \(S\), then
\[
\operatorname{SynState}_{\mathbb A,\mathbb B}(S)
\]
is obtained from
\[
\int_{\mathbb A}X
\]
by passing to a labelled internal subcategory and then to a labelled regular
quotient.

\item Relative to any behaviour-operation doctrine \(\mathcal O\), behaviour
varieties are classes of behaviour subobjects
\[
\mathcal V_{\mathbb A,\mathbb B}
\subseteq
\operatorname{Sub}_{A_0\times B_0}
(\mathsf{Beh}_{\mathbb A,\mathbb B}),
\]
not single subobjects of the final behaviour system.

\item If the state-size class
\[
\mathcal K_Z\subseteq\mathcal E/Z
\]
is closed under subobjects and is \(\mathcal O\)-closed, then stable classes
of closed \(\mathcal K\)-state recognizers correspond to
\(\mathcal K\)-recognizable \(\mathcal O\)-Mealy behaviour varieties by
\[
\mathcal P
\longmapsto
\left\{
S
\ \middle|\
\mathsf D(S)\in\mathcal P
\right\}
\]
and
\[
\mathcal V
\longmapsto
\left\{
D\in\mathsf{SRec}^{\mathcal K}_{\mathbb A,\mathbb B}
\ \middle|\
D\in\mathcal V
\right\}.
\]

\item For a category of state-admissible interface translations compatible
with the chosen state-size classes, the preceding state correspondence
globalizes objectwise and is compatible with inverse images of behaviour
specifications.

\item In the one-object Set case
\[
\mathbb A=M,
\qquad
\mathbb B=N,
\]
the residual closure of a specification
\[
S\subseteq\mathsf{Beh}_{M,N}
\]
is
\[
\mathsf D(S)
=
\{\partial_r\lambda\mid r\in M,\ \lambda\in S\}.
\]

\item If
\[
M=I^\ast,
\]
then
\[
\mathsf{Beh}_{I^\ast,N}\cong N^{I^+}
\]
and
\[
\mathsf D(S)
=
\{\partial_w\beta\mid w\in I^\ast,\ \beta\in S\}.
\]
For a generator \(i\in I\),
\[
\omega(i,\partial_w\beta)=\beta(wi),
\]
while for a composite word \(u=i_1\cdots i_n\), the folded output is
\[
\lambda_\beta(w,u)
=
\beta(wi_1)\beta(wi_1i_2)\cdots
\beta(wi_1\cdots i_n).
\]
\end{enumerate}
\end{theorem}

\begin{proof}
Parts (1)--(4) are
Definition~\ref{def:state-eilenberg-realized-behaviour-object},
Definition~\ref{def:state-eilenberg-recognition},
Proposition~\ref{prop:state-eilenberg-realized-behaviours-closed}, and
Corollary~\ref{cor:state-eilenberg-recognition-implies-closure-contained}.
Part (5) is
Theorem~\ref{thm:state-eilenberg-state-syntactic-divisibility}.
Parts (6) and (7) are
Definition~\ref{def:state-eilenberg-synstate} and
Corollary~\ref{cor:state-eilenberg-labelled-state-divisibility}.
Parts (8) and (9) are
Definition~\ref{def:state-eilenberg-behaviour-operation-doctrine},
Definition~\ref{def:state-eilenberg-local-behaviour-variety}, and
Theorem~\ref{thm:state-eilenberg-state-correspondence}.
Part (10) is
Theorem~\ref{thm:state-eilenberg-global-state-correspondence}.
Parts (11) and (12) are the one-object and free-input specializations
described in Sections~\ref{sec:state-eilenberg-one-object-set} and
\ref{sec:state-eilenberg-free-input-specialization}.
\end{proof}

The next chapter passes from the closed residual state recognizer
\[
D\hookrightarrow\mathsf{Beh}_{\mathbb A,\mathbb B}
\]
to the effective transition-output quotient determined by the labelled
input-discrete category
\[
\int_{\mathbb A}D\longrightarrow\mathbb A^{\mathrm{op}}.
\]
In the one-object case this becomes the quotient of the input monoid by the
syntactic transition-output congruence. This quotient-level construction is
the algebraic layer on which the Reiterman theory is formulated.